% ═══════════════════════════════════════════════════════════════
% MT-01b: Minitest Vorstellen + ser
% Spanisch Klasse 7 - Sequenz 1: Empezamos
% ═══════════════════════════════════════════════════════════════
% Dauer: 10 Minuten | Punkte: 28 BE
% ═══════════════════════════════════════════════════════════════

\documentclass[11pt, a4paper]{article}

\usepackage[utf8]{inputenc}
\usepackage[T1]{fontenc}
\usepackage[ngerman]{babel}
\usepackage{helvet}
\renewcommand{\familydefault}{\sfdefault}

\usepackage[margin=2cm]{geometry}
\usepackage{enumitem}
\usepackage{tabularx}
\usepackage{booktabs}
\usepackage{xcolor}
\usepackage{fancyhdr}
\usepackage{lastpage}

% FARBEN
\definecolor{spanisch}{HTML}{c41e3a}
\definecolor{grau}{HTML}{888888}
\definecolor{hellgrau}{HTML}{f5f5f5}

\pagestyle{fancy}
\fancyhf{}
\renewcommand{\headrulewidth}{0.4pt}
\renewcommand{\footrulewidth}{0.4pt}

\lhead{Spanisch -- Klasse 7}
\chead{\textbf{MT-01b}}
\rhead{\today}

\lfoot{10 Minuten}
\cfoot{}
\rfoot{Seite \thepage\ von \pageref{LastPage}}

\newcommand{\punktebox}[1]{\hfill\fbox{\small \_\_\_/#1}}
\newcommand{\luecke}[1]{\rule{#1}{0.4pt}}
\newcommand{\es}[1]{\textcolor{spanisch}{\textbf{#1}}}

\begin{document}

% ═══════════════════════════════════════════════════════════════
% KOPFBEREICH
% ═══════════════════════════════════════════════════════════════

\begin{center}
    {\Large\bfseries\textcolor{spanisch}{Minitest: Vorstellen + ser}}
\end{center}

\vspace{5pt}

\noindent
\begin{tabularx}{\textwidth}{|X|X|X|}
\hline
\textbf{Name:} \luecke{4cm} &
\textbf{Klasse:} 7 &
\textbf{Datum:} \luecke{2cm} \\
\hline
\end{tabularx}

\vspace{10pt}

\noindent
\fcolorbox{spanisch}{hellgrau}{%
\parbox{\dimexpr\textwidth-2\fboxsep-2\fboxrule}{%
\textbf{Zeit:} 10 Minuten | \textbf{Punkte:} 28 BE | \textbf{Hilfsmittel:} keine
}}

\vspace{15pt}

% ═══════════════════════════════════════════════════════════════
% AUFGABE 1: Konjugation ser (8 BE)
% ═══════════════════════════════════════════════════════════════

\noindent\textbf{\textcolor{spanisch}{Aufgabe 1 -- Das Verb ser konjugieren}} \punktebox{8}

\noindent Schreibe die richtige Form von \es{ser} (sein) neben das Personalpronomen.

\vspace{10pt}

\begin{tabularx}{\textwidth}{|l|X|}
\hline
\textbf{Pronomen} & \textbf{ser} \\
\hline
yo & \luecke{3cm} \\[0.6em]
\hline
tú & \luecke{3cm} \\[0.6em]
\hline
él / ella & \luecke{3cm} \\[0.6em]
\hline
nosotros & \luecke{3cm} \\[0.6em]
\hline
\end{tabularx}

\vspace{15pt}

% ═══════════════════════════════════════════════════════════════
% AUFGABE 2: Lückentext ser (8 BE)
% ═══════════════════════════════════════════════════════════════

\noindent\textbf{\textcolor{spanisch}{Aufgabe 2 -- ser einsetzen}} \punktebox{8}

\noindent Ergänze die richtige Form von \es{ser}.

\vspace{10pt}

\begin{enumerate}[label=\alph*)]
    \item Yo \luecke{2.5cm} de Alemania.
    \item Tú \luecke{2.5cm} estudiante.
    \item Ella \luecke{2.5cm} de España.
    \item Nosotros \luecke{2.5cm} de Rostock.
\end{enumerate}

\vspace{15pt}

% ═══════════════════════════════════════════════════════════════
% AUFGABE 3: Redemittel (6 BE)
% ═══════════════════════════════════════════════════════════════

\noindent\textbf{\textcolor{spanisch}{Aufgabe 3 -- Redemittel zuordnen}} \punktebox{6}

\noindent Verbinde die spanischen Ausdrücke mit der deutschen Bedeutung. Schreibe den Buchstaben in die Lücke.

\vspace{10pt}

\begin{minipage}[t]{0.55\textwidth}
\begin{enumerate}[label=\arabic*.]
    \item \es{Me llamo...} \luecke{1cm}
    \item \es{Soy de...} \luecke{1cm}
    \item \es{¿Cómo te llamas?} \luecke{1cm}
    \item \es{¿De dónde eres?} \luecke{1cm}
\end{enumerate}
\end{minipage}
\hfill
\begin{minipage}[t]{0.4\textwidth}
\begin{itemize}[label=]
    \item a) Ich heiße...
    \item b) Ich bin aus...
    \item c) Wie heißt du?
    \item d) Woher bist du?
\end{itemize}
\end{minipage}

\vspace{15pt}

% ═══════════════════════════════════════════════════════════════
% AUFGABE 4: Dialog (6 BE)
% ═══════════════════════════════════════════════════════════════

\noindent\textbf{\textcolor{spanisch}{Aufgabe 4 -- Dialog vervollständigen}} \punktebox{6}

\noindent Ergänze die Antworten von Person B sinnvoll auf Spanisch.

\vspace{10pt}

\noindent
\textbf{A:} ¡Hola! ¿Cómo te llamas?

\vspace{0.6em}
\noindent
\textbf{B:} \luecke{8cm}

\vspace{0.8em}
\noindent
\textbf{A:} ¿De dónde eres?

\vspace{0.6em}
\noindent
\textbf{B:} \luecke{8cm}

\vspace{0.8em}
\noindent
\textbf{A:} ¡Mucho gusto!

\vspace{0.6em}
\noindent
\textbf{B:} \luecke{8cm}

% ═══════════════════════════════════════════════════════════════
% PUNKTEÜBERSICHT
% ═══════════════════════════════════════════════════════════════

\vspace{20pt}
\hrule
\vspace{10pt}

\noindent
\begin{tabularx}{\textwidth}{|X|c|c|c|c||c|}
\hline
\textbf{Aufgabe} & 1 & 2 & 3 & 4 & \textbf{Gesamt} \\
\hline
\textbf{max. Punkte} & 8 & 8 & 6 & 6 & \textbf{28} \\
\hline
\textbf{erreicht} & & & & & \\[0.6em]
\hline
\end{tabularx}

\vspace{10pt}

\begin{flushright}
\textbf{Note:} \luecke{2cm}
\end{flushright}

\end{document}
